\chapter*{Lời Mở Đầu}
\addcontentsline{toc}{chapter}{Lời Mở Đầu}
\markboth{Lời Mở Đầu}{Lời Mở Đầu}

\begin{truyen}{Lời Mở Đầu}{Opening Words}
\chuhoa{L}{ớp học là một nơi rất lạ: nhiều điều nhỏ xảy ra mỗi ngày, nhưng chính những điều nhỏ ấy lại ở lại rất lâu trong trí nhớ con người.}
\textit{(The classroom is a strange place: many small things happen every day, yet those small things remain in memory for a very long time.)}

Đó có thể là một câu học sinh nói rất khẽ, một lỗi sai khiến cả lớp bật cười, một lần giơ tay đầy run, hay một khoảnh khắc em nào đó hiểu bài và sáng hẳn mắt lên.
\textit{(It may be a sentence a student says softly, a mistake that makes the whole class laugh, a trembling raised hand, or a moment when someone finally understands and their eyes brighten.)}

Người ngoài có thể nghĩ đó chỉ là chuyện vụn vặt của một giờ học.
\textit{(An outsider may think these are just tiny incidents from an ordinary class.)}

Nhưng thật ra, chính từ những điều vụn ấy mà học sinh học cách nghĩ, cách hỏi, cách sai, cách sửa, cách cố gắng và cả cách lớn lên.
\textit{(But in truth, it is from these small things that students learn how to think, ask, be wrong, correct themselves, keep trying, and grow up.)}

Quyển này đi vào lớp học từ chính những mẩu nhỏ như thế. Không phải để làm mọi thứ trở nên to tát, mà để giữ lại những bài học rất thật thường bị trôi qua quá nhanh.
\textit{(This volume enters the classroom through such small fragments. Not to make everything seem grand, but to keep the very real lessons that usually pass too quickly.)}

Nếu em là học sinh, mong em thấy mình ở đâu đó trong những trang này.
\textit{(If you are a student, I hope you see yourself somewhere in these pages.)}

Nếu em là giáo viên, mong em thấy lại vì sao lớp học không chỉ là nơi truyền bài, mà còn là nơi từng con người đang thành hình.
\textit{(If you are a teacher, I hope you remember why the classroom is not only where lessons are delivered, but where human beings are being formed.)}
\end{truyen}

\begin{baihoc}
Trong lớp học, điều nhỏ không hề nhỏ. Nó thường là nơi bài học thật sự bắt đầu.
\textit{(In a classroom, small things are not small. They are often where the real lesson begins.)}
\end{baihoc}

\begin{ghinhoanh}
\textbf{Short English to remember:}

Small classroom moments can carry lasting lessons.

\end{ghinhoanh}

\begin{tuhocnhanh}
1. Vì sao những điều rất nhỏ trong lớp học lại đáng nhớ? \textit{Why can very small classroom moments become memorable?}

2. Một giờ học tốt thường được tạo nên từ những điều gì? \textit{What kinds of things usually make a good lesson?}

3. Em muốn giữ lại bài học nhỏ nào từ lớp học của chính mình? \textit{What small lesson from your own classroom would you keep?}
\end{tuhocnhanh}

\ngancach